\section{Problem Setting and Method}
\label{sec:method}

\paragraph{Leaderboard and rules.}
A leaderboard is a score matrix $x\in[0,100]^{M\times n}$ of $M$ models on $n$ tasks (sub-datasets).
An aggregation rule $R$ maps a matrix to a ranking of the models.
We study score rules, which compute a scalar per model: the arithmetic \emph{mean} (the equal-task mean; on the Open LLM Leaderboard, the mean of the six official normalised task scores is the official average, while the raw-accuracy columns are a separate variant; we do not claim that it reproduces HELM's own aggregate), the \emph{median}, the \emph{minimum}, and the \emph{geometric mean}; and two ordinal rules, which use only within-task order: the \emph{Borda} count and the requirement to be a strict \emph{Condorcet winner} (beating every rival in pairwise comparison, where a task won counts 1, a task tied counts $\tfrac12$, and the target wins if its total exceeds $n/2$; with ties this is not the same as winning more than $n/2$ tasks strictly).

\paragraph{Intervention and investment.}
Fix a target model $i$. The developer may raise its scores to any $y_t\in[x_{it},c_t]$, where $c_t$ is a cap on task $t$; all rivals' raw scores stay fixed, there are no spillovers to other tasks, and no rival responds.
After the edit, ordinal aggregates of rivals are recomputed because the target may pass them.
The target is \emph{first} if it beats every rival under $R$ by a margin $\varepsilon$ (score rules; $\varepsilon=10^{-3}$ points unless stated) or strictly (ordinal rules).
Strict victory over a continuous score can have an infimum that is not attained, so $\varepsilon$ defines the margin version we report, and we report sensitivity to it; a target that already wins has cost zero.

\emph{Investment is a score counterfactual, not a training cost.}
We measure it by three separate optima, each its own optimisation problem: $L_1^\ast$, the least total number of points added; $K^\ast$, the fewest tasks touched; and $L_\infty^\ast$, the smallest largest single-task lift when any subset may be raised.
Where a solver minimises one of them, the others are only reported descriptors at its solution.
Raising all tasks by a common amount is a different quantity and is not called $L_\infty^\ast$.
A target is \emph{infeasible} if no allowed target row under the cap satisfies the winning condition (a cap that binds at the optimum does not make a target infeasible); we report the feasible fraction separately from costs over feasible targets, never averaging an infinite cost, together with the smallest cap at which the target becomes feasible.
For the Condorcet requirement, the cost is that of a sufficient condition and therefore an upper bound for any particular Condorcet-consistent rule (Schulze, Kemeny, \dots); it is not that rule's exact cost.

\paragraph{Caps, scales and cohort (modelling assumptions).}
The cap is the central assumption and the headline result is investment as a function of it.
We use a loose ceiling ($c_t=100$), a frontier cap (best observed score on $t$ plus a margin, never below the current score), and an empirical cap (the current score plus a high percentile of observed fine-tune-minus-base score gains on $t$, never below the current score and never above the ceiling); further caps appear in the cap sweep.
Frontier and empirical caps are assumptions, not evidence that such scores are attainable; empirical caps carry survivorship bias and use one gain quantile for all targets.
Open LLM Leaderboard scores are on unrelated scales, so raw scores, per-task standardised scores and rank percentiles are distinct variants (Section~\ref{sec:exp}).
Geometric means use only positive-domain scales; $z$-scores are used only for mean, median, minimum and pooled $L_1$.
Transform parameters are fixed from the original snapshot and edits and caps are mapped through the same transform.
In every comparison across rules the targets, rivals and caps are identical (a fixed cohort); where the ordinal solvers are too slow, a rank-stratified sample is used for \emph{all} rules in that table and results are reported by rank bin.
Task bootstraps and evaluation-noise simulations are sensitivity analyses conditional on these assumptions and do not validate them.

\paragraph{Closed forms for the score rules.}
Let $\tau$ be the best rival's value under the rule, $b=\tau+\varepsilon$, $h_t=\max(c_t-x_{it},0)$ the headroom, and $x=x_{i\cdot}$.
Write $d^\ast$ for the smallest $d$ such that $y=\min(x+d,c)$ makes the target first: since every rule is monotone in the target's own scores, any set of edits that wins can be replaced by this uniform-lift edit with no larger largest lift, so $L_\infty^\ast=d^\ast$ for every rule (Proposition~4).
\begin{enumerate}\setlength\itemsep{0pt}
\item \textbf{Mean.} Let $\mathrm{need}=nb-\sum_t x_t$. The target is already first if $\mathrm{need}\le 0$; it is feasible iff $\sum_t h_t\ge\mathrm{need}$; then $L_1^\ast=\mathrm{need}$ and $K^\ast$ is the fewest tasks whose largest headrooms sum to at least $\mathrm{need}$.
\item \textbf{Median, $n$ odd.} The target needs at least $m=(n+1)/2$ tasks with $y_t\ge b$. With $a=\#\{t:x_t\ge b\}$, $K^\ast=\max(0,m-a)$; $L_1^\ast$ is the sum of the $K^\ast$ smallest values of $b-x_t$ over tasks below $b$ with $c_t\ge b$; it is infeasible iff fewer than $K^\ast$ such tasks exist. (For even $n$, as on the six-task Open LLM Leaderboard, the median is the mean of two middle values and the solver is used; it is checked against brute force.)
\item \textbf{Minimum.} $K^\ast=\#\{t:x_t<b\}$, $L_1^\ast=\sum_t(b-x_t)^+$, feasible iff $c_t\ge b$ on every task below $b$.
\end{enumerate}
These formulas explain the cap dependence we find: a rule that requires every task to reach $b$ (minimum) or a majority of tasks (median) is blocked as soon as the cap lies below $b$ on those tasks, whereas the mean can trade headroom across tasks.
They are conditional on the model above; they are not statements about training.
We verify the formulas, and Proposition~4 for the geometric mean, Borda and Condorcet as well, against the exact solvers on task subsets, targets and caps sampled from the two matrices (31{,}080 checks, all agreeing; Appendix, Table~\ref{tab:theory}); the proofs are short and are given there.
